\scanpage{36}
N.B. Fixons un pt $y_0$ de $Y$. Alors le groupe
des familles algébriques $Y(k)\unsure{-|\longrightarrow}A^i_{\mathrm{alb}}(X)$ est
isom à $A^i_{\mathrm{alb}}(X)\times\operatorname{Hom\,alg}(A_Y,A^i_{\mathrm{alb}}(X))$, où le
deuxième facteur désigne les familles algébriques
qui sont des hom de groupes, i.e. nulles à
l'origine. On peut donc \deleted{se} limiter à étudier
les hom ``alg.'' $A_{\deleted{\unsure{Y}}}(k)\longrightarrow A^i_{\mathrm{alb}}(X)$, \inserted{où $A$ est une VA,} et interprétant
$A$ comme $A_{0\,\mathrm{alb}}(A)=A^\alpha_{\mathrm{alb}}(A)^{(\text{alg.équiv.}\,0)}$ $(\alpha=\dim A)$
on peut considérer \deleted{\unsure{le}} ces objets comme des
cas particuliers des \uline{hom algébriques}
\margin{$A(k)\simeq P^{\unsure{\alpha|n}}(A)\longrightarrow P^i(X)$}
\[
\varphi:P^j(X')\longrightarrow P^i(X).
\]

Un hom de gpes $\varphi$ est dit ``algébrique''
s'il existe un entier $N>0$ tel que
$N\varphi$ soit définissable par un cycle algébrique sur $X\times\unsure{Y'|X'}$
(donc $\varphi$ serait définissable par un cycle alg.
à coeff. \uline{rationnels} --- mais attention, un cycle
à coeff. rationnels \uline{quelconque} ne définit pas en
général un hom $P^j(X')\longrightarrow P^i(X)$, $=$ si les
degrés sont les bons ....). Ces hom \uline{se composent}
% archive sheet number: 35
